\scanpage{31}

% L1–2
Il résulte de (3) et (2) appliqué à $i=\omega$ que l'on a
\[
 \alpha_\omega(d)\subset d_\omega\subset\alpha_\omega(d)\cup
 \bigcup_{\substack{\deleted{\unsure{i<\omega}}\\(S,\lambda)\in\sigma}}
 \ \bigcup_{u\in\underline{\mathrm{Hom}}_\lambda(S,C_\omega)}
 \qquad\substack{\text{mod isom de}\\\text{cônes}}
\]
% L3–5
et pour prouver $3^\circ)$ il faut donc prouver
\[
 \begin{aligned}
 &\forall(S,\lambda)\in\sigma\text{ et }u\in\underline{\mathrm{Hom}}_\lambda(S,C_\omega),
 \ \deleted{\text{exi}}\exists\,i<\omega\text{ et un}\\
 &\text{hom }u_i:S\to C_i\text{ tel que }u_i(\lambda)\subset d_i\ \inserted{\text{(mod iso de cônes)}}
 \text{ et }u\simeq\tau_{\omega,i}\circ u_i
 \end{aligned}\tag*{(\deleted{\unsure{3 bis}}\inserted{4})}
\]

% L6–7
\deleted{Cela résulte facilement du fait que tt se ramène à un $C_i$, de d) et du fait que}
Enfin, \deleted{\unsure{pas}} examinons $1^\circ)$. \deleted{A tout \unsure{prendre}}

% L8–10
$a_1)$ Sera conséquence de d) si nous notons que tout foncteur $u:S-C_i$ se factorise par un $C_i$ -- pourvu qu'on prenne $\pi\geq\operatorname{Card}\operatorname{Fl}S$ pour tt $(S,\lambda)\in\sigma$.

% L11–12
$a_2)$ Sera conséquence de e), utilisant B (les $(R_i,\delta_i)\in\sigma$)) et \deleted{\unsure{(3 4 i)|(4)}} et (3) (pour $i+1$ au lieu de $i$).

% L13–14, crossed block
\deleted{Donc il reste à \unsure{donc|donner} dégager des conditions sur la construction \unsure{$\alpha|K$} qui assurent $(\unsure{4|3})$, et d' \inserted{enfin à} expliciter une construction $K$.}

% L15, left
\margin{peut remonter !}

% L15–18
Prouvons 4). On sait que $u:S-C_\omega$ se factorise par un \inserted{$u_i:S\to$} $C_i$ (pourvu qu'on prenne $\pi\geq\operatorname{card}\operatorname{Fl}S$ $\forall(S,\lambda)\in\sigma$), donc quitte à composer avec $\tau_i$, on peut supposer que \unsure{cel|cela|ceci} se prolonge en $\bar u_i:\bar S-C_i$ tel que $\bar u(\bar\lambda)\in d_i$, donc $\tau_{\omega,i}\circ\bar u_i$ est un hom $\bar S\to C$ qui est $\bar\lambda$ exact. \inserted{grâce à (3).} \deleted{Il est donc} Comme
% L19–24
$u$ est $\lambda$-exact, $\bar u(i_c)$ est un isom pour tt $c\in\lambda$, donc $\bar u_i(i_c)$ devient un isom pour tt $c\in\lambda$, donc les $\tau_{ji}\bar u_i(i_c)=\bar u_j(i_c)$ ont la $=$ propriété pour $j\in T_\omega$ assez grand (pourvu que $\pi\geq\operatorname{card}\lambda=\operatorname{card}\{c\}$) donc \deleted{\illegible{1}} les \deleted{éléments} cônes de $u_j(\lambda)$ sont \uline{isom} à ceux de $\bar u_j(\bar\lambda)$, comme ceux-ci sont dans $d_i$, il en est de $=$ de ceux là à isom près, cqfd.
% archive sheet number: 30
